
@@M0@@
devient ainsi un module sur l'anneau @@M1@@ (et même un module gradué sur l'anneau gradué @@M2@@, à condition de changer de signe de la graduation de @@M3@@ en posant @@M4@@ @@M5@@).
Considérons maintenant un morphisme propre
@@M6@@
de schémas noethériens. Alors @@M7@@ induit un homomorphisme
@@M8@@
d'anneaux gradués (V), et d'après 1.1.8 un homomorphisme
@@M9@@
de groupes abéliens gradués. Alors la formule de projection (IV 2.11) et (3.1.1) donnent par passage aux gradués associés la formule de projection graduée
@@M10@@

\subsection*{3.2.}
Soit @@M11@@ un morphisme projectif d'intersection complète (Exp. VIII) et de dimension virtuelle @@M12@@ de schémas noethériens, @@M13@@ ayant un @@M14@@-module inversible ample. Dans Exp.VIII on a démontré que @@M15@@ induit un homomorphisme
@@M16@@
qui satisfait d'ailleurs à la formule de projection (Exp. IV) :
@@M17@@
